%% !TEX root = ../manual_gaolv5.tex
% Chapter 9.
%% Chapter 9: Input and output. Sources: gaol.tex of f5bf180 lines 4904-5531
%% (reading intervals, the input format, writing intervals, converting to
%% strings, the output format, the number of digits, the example); the brief
%% named no doc/ page for this chapter.
%%
%% Copyright (c) 2001 Swiss Federal Institute of Technology, Switzerland
%% Copyright (c) 2002-2016 Laboratoire d'Informatique de
%%                         Nantes-Atlantique, France
%% Copyright (c) 2017-2020 Laboratoire des Sciences du Numérique de Nantes,
%%                         France
%% Copyright (c) 2026 ENSTA, France
%%
%% By: Frederic Goualard, Jordan NININ
%%
%% GAOL is a software distributed WITHOUT ANY WARRANTY. Read the associated
%% COPYING.LIB file for information.
\chapter{Input and output}\label{chap:io}

An interval enters a program and leaves it as text: a literal in a source
file, a line of a data file, a string typed by the user, a result printed on
the screen. A decimal number is seldom a double, and a double is seldom
written exactly with a few decimal digits. Reading and writing must therefore
preserve the containment property (\pxref{sec:containment}): the text read
gives an interval that contains the number or the set it writes, and the text
written for an interval in the formats \element{bounds} and \element{hexa}, or
by \ieee{intervalToText}, read back, gives an interval that contains it. GAOL
v5 does both, whatever the C library and the locale of the
program.\newinvfive

\medskip
In this chapter, we describe how GAOL v5 reads and writes intervals. The
first section covers the reading: the operator \code{>{}>} of the input
streams, then the language of the strings that it and \code{textToInterval()}
accept, with its numbers, the uncertain form of IEEE 1788-2015, its functions,
its bounds and infinities, the limits of the reader and its errors. The second
section covers the writing: the operator \code{<{}<}, the exact text of
\code{exact\_string()}, the conversion to \code{std::string}, the five output
formats, the number of digits, and a complete program. The functions that
read and write with the names of the standard, \ieee{textToInterval},
\ieee{intervalToText}, \ieee{intervalToExact} and \ieee{exactToInterval}, belong to the namespace
\code{gaol\_ieee1788} and are described in Chapter~\ref{chap:ieee1788}
(\pxref{sec:ieee1788-text}).

\section{Reading intervals}\label{sec:reading-intervals}

An interval is read from a stream with the operator \code{>{}>}, as a number
is, and built from a string with \code{textToInterval()}
(Chapter~\ref{chap:creation}). Both read the language of
Section~\ref{sec:input-format}.
\cindex{interval!reading}%

\begin{deffun}{istream\&}{operator>{}>}{(istream\& \var{in}, interval\& \var{I})}
  Reads a line from the input stream \var{in}, and assigns to \var{I} the
  interval it writes. The blanks before the line, the ends of lines
  included, are skipped first, as before a number, even under
  \code{std::noskipws}, which keeps a number from skipping them: an empty
  line, or one holding only blanks, is never a line to read. An interval
  cannot be written on several lines.

  If the line is syntactically ill-formed, \var{I} is set to the empty set.
  A \code{std::invalid\_argument} exception is then thrown
  (\pxref{sec:exceptions}) if the library was compiled with exceptions
  enabled (\pxref{sec:install-developers}); if they were disabled, an error
  message is printed to \code{cerr} and the program aborts. An ill-formed
  line also sets the \code{failbit} of \var{in}, which \code{in.clear()}
  resets before reading on. A second argument of \code{nth\_root} that is not an integer, as in
  \code{nth\_root(8,1.5)}, throws \code{std::domain\_error} instead; a call with
  a wrong number of arguments throws \code{std::invalid\_argument}, and an
  argument outside the domain of a function gives the empty set, as
  \code{sqrt(-1)} does.

  Where no line is left, \var{in} gets its \code{failbit}, \var{I} is left
  as it was, and nothing is thrown, as for a number: \code{while (in >{}> x)}
  reads the lines one after the other, the blank ones skipped, and ends with
  the input, a file ending with an empty line included.\newinvfive{} GAOL~4
  read the empty text left, and threw \code{input\_format\_error}, as it did
  on an empty line, so that \code{in >{}> d >{}> x} over the lines
  \code{1.5} and \code{[1, 2]} threw at the end of the first one.
\begin{example}
std::istringstream in("[1, 2]\n[0.1, 0.2]\n");
interval a, b;
in >> a >> b;
cout << a << ' ' << b;
@outputs^[1, 2] [0.09999999999999999, 0.2000000000000001]~
\end{example}
\begin{example}
std::istringstream in("1.5\n\n[1, 2]\n\n");
double d;
interval x;
in >> d >> x;
cout << d << ' ' << x << ' ';
cout << (in >> x ? "more" : "end");
@outputs^1.5 [1, 2] end~
\end{example}
\end{deffun}

\subsection{The input format}\label{sec:input-format}

A string to be translated into an interval has the syntax below, with the
terminals in lower case and the non-terminals in italic upper case.
\cindex{input format}%

\begin{Verbatim}[commandchars=\@\^\~,fontsize=\small]
@NT^expr~
	: @NT^number~
	| @NT^uncertain_number~ @rem^@hfill// An interval literal, never a bound by itself~
	| @NT^literal~
	| @TM^dmin~  @rem^@hfill// Smallest positive normal floating-point number~
	| @TM^dmax~ @rem^@hfill// Largest positive floating-point number~
	| @TM^pi~
	| @TM^inf~ @rem^@hfill// [MAX, +infinity] as an expression, +infinity as a bound~
	| @TM^infinity~ @rem^@hfill// The same~
	| @TM^empty~ @rem^@hfill// Empty interval~
	| @NT^expr~ @TM^+~ @NT^expr~
	| @NT^expr~ @TM^-~ @NT^expr~
	| @NT^expr~ @TM^*~ @NT^expr~
	| @NT^expr~ @TM^/~ @NT^expr~
	| @TM^-~ @NT^expr~
	| @TM^+~ @NT^expr~
	| @NT^function_call~
	| @TM^(~ @NT^expr~ @TM^)~
	;

@NT^literal~
	: [ @NT^expr~ ]
	| [ @NT^expr~ @TM^,~ @NT^expr~ ]
	| [ @NT^expr~ @TM^,~ ] @rem^@hfill// Right bound +@ensuremath^@infty~~
	| [ @TM^,~ @NT^expr~ ] @rem^@hfill// Left bound -@ensuremath^@infty~~
	| [ @TM^,~ ] @rem^@hfill// Entire real line~
	| [ @TM^entire~ ] @rem^@hfill// Entire real line~
	| [ ] @rem^@hfill// Empty interval~
	| < @NT^expr~ @TM^,~ @NT^expr~ >
	;

@NT^function_call~
	: @NT^name~ @TM^(~ @NT^expr~ @TM^)~
	| @NT^name~ @TM^(~ @NT^expr~ @TM^,~ @NT^expr~ @TM^)~
	| @NT^name~ @TM^(~ @NT^expr~ @TM^,~ @NT^expr~ @TM^,~ @NT^expr~ @TM^)~
	;
\end{Verbatim}

\noindent
Every string is an expression, computed with interval arithmetic. The
operators \TM{*} and \TM{/} bind tighter than \TM{+} and \TM{-}, the four
associate to the left, and a sign binds tighter than all of them:
\code{"1-2-3"} is $-4$ and \code{"2+3*4"} is 14. The
literals of intervals, the numbers, the constants and the functions combine
with the operators in any order, and a literal stands wherever a number may,
in a bound too.\newinvfive{} \code{1+[1,2]}, \code{2*cos([0,1])} and
\code{[cos([0,1]), 2]} are read. \code{[empty]} is the form
\code{[\NT{expr}]} of the expression \code{empty}.

\begin{Rem}
GAOL~4 read a string with two grammars, one for the intervals between
brackets and one for the numbers. It refused \code{1+[1,2]} and
\code{[1,2]+1+[1,2]}, while it read \code{[1,2]*2}. GAOL v5 reads every
string with the one grammar above.\newinvfive
\end{Rem}

Spaces, tabs and newlines are not significant except in numbers, and the
case of letters is ignored, as IEEE 1788-2015 has it for its literals:
\code{[Empty]}, \code{[1, Inf]}, \code{SIN(1)} and \code{1E3} are read,
whatever the locale of the program.\newinvfive{} The sign ``$+$'' before a
number and before \code{inf} is optional.

\paragraph{Numbers.} A \NT{number} is written in decimal, with an optional
point and an optional exponent (\code{12}, \code{0.1}, \code{1e-3}), in
hexadecimal (\code{0x1f}, \code{0x1.3p-1}), in octal (\code{0o17}) or in
binary (\code{0b101}, \code{0b1.1p3}); the digits after \code{0x}, \code{0o}
or \code{0b} have an optional point and an optional binary exponent after a
\code{p}.\newinvfive{} A number that begins with 0 is decimal: \code{010} is
10, as IEEE 1788-2015 has it (9.7.2), and octal is written \code{0o10}, as
in Python.
\cindex{number!reading}%

Numbers are read exactly, whatever their number of digits: a number that is
a double is read as that double, and any other one is replaced by the
tightest interval of doubles containing it.\newinvfive{} GAOL~4 read them
with the \code{strtod()} of the C library, which did not enclose \code{0.1}
where it rounds to nearest in every rounding direction (the C runtime of
Windows, musl on 64-bit ARM processors). The point of a number is a point
whatever the locale of the program, a locale writing a decimal comma
included (as \code{setlocale(LC\_ALL, "")} sets under a French
one).\newinvfive{} GAOL~4 read \code{0.1} as 0 there, \code{strtod()}
stopping at the point. The modes of the processor that flush the subnormal
numbers to zero, flush-to-zero and denormals-are-zero, which a program
linked with \option{-Ofast} has (\pxref{sec:ftz}), do not change what a
number is read as: \code{1e-310} is read as the tightest interval of doubles
containing it, and \code{0} as 0.\newinvfive{} Reading a long number takes a
time quadratic in its number of digits: \code{"0."} followed by $n$ digits
\code{3} takes 0.011~s for $n = 20\,000$, 0.18~s for 80\,000 and 2.8~s for
320\,000 (about 28~s expected for a million), and the uncertain form 2.2 to 2.4
times as long (Intel i7-1185G7, GCC~15.2, the best of five runs,
7 November 2026).\newinvfive

\paragraph{The uncertain form.} An \NT{uncertain\_number} is written as
IEEE 1788-2015 defines it (9.7.4 and 12.11.3).
\cindex{uncertain form}%

\begin{Def}
An \dfn{uncertain number} is a decimal number followed by \code{?}, a radius
in units of its last digit (half a unit when it is left out, and an infinite
radius for a second \code{?}), an optional direction, \code{u} or \code{d},
which keeps only the radius above or below, and an optional exponent.
\end{Def}

\noindent
\code{3.56?1} is \itv{3.55}{3.57}, \code{3.56?1e2} is \itv{355}{357}, and
\code{-10??u} is \itv{-10}{+\infty}.\newinvfive{} An uncertain number is an
interval literal of its own, not a bound by itself: \code{"[5?1]"} and
\code{"[1, 3.56?1]"} throw \code{std::invalid\_argument}, as IEEE 1788-2015
has it (12.11.4), while \code{"3.56?1*2"} is \itv{7.1}{7.14}, and a bound
computed from one is read, \code{"[2*5?1]"} being \itv{8}{12}. A sign written
next to an uncertain number, as \code{10-3.56?1} or \code{10+3.56?1}, is
read as a part of the number, and the text throws
\code{std::invalid\_argument}, as does \code{10 -3.56?1}: a blank is needed
after the sign, \mbox{\code{10 - 3.56?1}}.\newinvfive

\paragraph{Functions.} A \NT{name} is the name of a function of GAOL, with
its arguments, whatever the case of its letters. The names of GAOL are read
by \code{gaol::textToInterval()}, from one string or from two, and by
\code{operator>{}>}; the names of IEEE 1788-2015 are read by
\code{gaol\_ieee1788::textToInterval()} (\pxref{sec:ieee1788-text}), as
each namespace has its own \code{pow}.\newinvfive{} The names of GAOL are
the following.

\begin{center}\small
\begin{tabular}{L{11cm}}
\code{sqr}, \code{sqrt}, \code{cbrt}, \code{pow}, \code{nth\_root},
\code{exp}, \code{exp2}, \code{exp10}, \code{expm1}, \code{exp2m1},
\code{exp10m1}, \code{log}, \code{log2}, \code{log10}, \code{log1p},
\code{log2p1}, \code{log10p1}, \code{sin}, \code{cos}, \code{tan},
\code{asin}, \code{acos}, \code{atan}, \code{atan2}, \code{sinh},
\code{cosh}, \code{tanh}, \code{asinh}, \code{acosh}, \code{atanh},
\code{sinpi}, \code{cospi}, \code{tanpi}, \code{asinpi}, \code{acospi},
\code{atanpi}, \code{atan2pi}, \code{hypot}, \code{compoundm1},
\code{rsqrt}, \code{sign}, \code{ceil}, \code{floor}, \code{trunc},
\code{integer}, \code{round\_ties\_to\_even}, \code{round\_ties\_to\_away},
\code{abs}, \code{min}, \code{max}, \code{inverse}, \code{erf}, \code{erfc},
\code{cancel\_minus}, \code{cancel\_plus}, \code{hull}, \code{intersect},
\code{fma}\\
\end{tabular}
\end{center}

\noindent
Each is the function of GAOL of that name (Chapters~\ref{chap:arithmetic}
and~\ref{chap:functions}); \code{pow} is the \code{pow} of the namespace
\code{gaol} (\pxref{sec:powers}), and \code{cbrt(x)} is
\code{nth\_root(x,3)}. A name that is not in the list, or a call with a
wrong number of arguments, throws \code{std::invalid\_argument}. The second
argument of \code{nth\_root} shall be a point interval that can be evaluated
as an integer, negative ones included; otherwise, a \code{std::domain\_error}
exception is thrown, or an error is reported, depending on whether the
exceptions were enabled at configuration time
(\pxref{sec:install-developers}).\newinvfive{} The integer may be of any size a
double holds: \code{"nth\_root(8,3000000000)"} is
\code{nth\_root(8, 3000000000LL)}, where it was refused beyond the
\code{int}s. One of $2^{63}$ or more in magnitude is taken as
$\pm(2^{63}-2)$: the integer powers are the same, and the enclosure of the
roots of that order holds those of the larger orders. The second argument of
\code{pown} and \code{rootn}, with the names of the standard, is read in the
same way. GAOL~4 read the trigonometric and hyperbolic functions,
\code{exp}, \code{log}, \code{pow}, \code{sqrt} and \code{nth\_root} only.

\paragraph{Bounds.} The notations ``$n$'' and ``[$n$]'' are equivalent, but for
\code{inf} and \code{-inf}, whose literals \code{[inf]} and \code{[-inf]} are
empty (see Infinities below), and for an uncertain number, which is no bound. The
expressions in the bounds are evaluated with interval arithmetic, intervals
included; the left (respectively the right) bound of the result is then
used, depending on the side it appears in: \code{"[cos([0,1]), 2]"} is
\itv{\cos 1}{2}. A bound left out is infinite.

An interval literal \code{[l, u]} whose left bound is above its right bound
is the empty set, as in IEEE 1788-2015: \code{"[2,1]"}, and also
\code{"[9007199254740993, 9007199254740992]"} and
\code{"[1.00000000000000001, 1]"}, whose numbers are compared exactly with
the doubles around them, a number that is not a double lying strictly
between the two doubles around it.\newinvfive{} GAOL~4 read the last two as
the points $2^{53}$ and 1. Two numbers that are not doubles and lie between
the same two doubles are compared exactly too:
\code{"[1.00000000000000002, 1.00000000000000001]"} is the empty set, and
\code{"[1.00000000000000001, 1.00000000000000002]"} is the interval from
the double below to the double above them, as is the same number written
twice. Three kinds of bounds are not told apart: \code{pi} and another
number lying between the same two doubles; a number far beyond the doubles
(about $10^{311}$ or more, or $10^{-331}$ or less, in decimal, about
$2^{1026}$ or more, or $2^{-1081}$ or less, in the other bases); and bounds
that are computed, as \code{1/3}. The literal is then the interval from the
lower bound of the first to the upper bound of the second (case (b) of
12.12.7, which GAOL signals by no exception: \pxref{sec:ieee1788-missing}).

The two bounds in the string \code{<a, b>} must be expressions that
evaluate to the same double whatever the rounding direction; a
\code{std::invalid\_argument} exception is raised otherwise, as for
\code{"<0.1, 0.1>"}, a number that is not a double being no bound of
\code{<a, b>}.\newinvfive{} GAOL~4 read
\code{"<0.1, 0x1.9999999999999p-4>"}, the double just below $0.1$, as the
point $[$\code{0x1.9999999999999p-4}$]$. A bound that is computed is a
double when its value is a point, as \code{0.5}, \code{(0.25+0.25)} and
\code{[0.5]} are; \code{"<[0.1], 0x1.9999999999999p-4>"} and
\code{"<0x1.5555555555556p-2, 1/3>"}, whose first or second bound has the two
doubles around the number for its value, throw \code{std::invalid\_argument}
too, where GAOL~4 read them as the point of the other bound.\newinvfive{}
The two doubles are compared by their bits, $-0$ and $+0$ being the same
bound, so that the mode of the processor that reads the subnormal numbers
as zero (denormals-are-zero) changes nothing: \code{"<0x1p-1074, 0x1p-1073>"},
two subnormal doubles, throws under it too, where a comparison of the bounds as
doubles would take them for one.\newinvfive

\begin{Rem}
GAOL~4 wrote the degenerate intervals in the form \code{<a, b>}, and the
texts it wrote for those whose digits are exact, as \code{<4, 4>}, are still
read. \code{operator<{}<} writes such an interval \code{[a]}, the literal of
IEEE 1788-2015 for a point, and any other with its two bounds
(\pxref{sec:output-format}).
\end{Rem}

\paragraph{Infinities.} An interval literal whose left bound is $+\infty$ or
whose right bound is $-\infty$ is the empty set, as \ieee{numsToInterval} of
IEEE 1788-2015 has it: \code{"[inf]"} and \code{"[-inf, -inf]"} are
empty.\newinvfive{} The expression \code{"inf"} alone is still
\itv{\text{MAX}}{+\infty}, and \code{"-inf"} \itv{-\infty}{-\text{MAX}}
(\pxref{par:inf_inf_note}).

\paragraph{Threads.} Several threads may read strings at once, in parallel:
the lexer is reentrant and the parser pure, and each reading of a string has
a state of its own.\newinvfive{} GAOL~4 kept that state in global variables,
and two threads reading a string each crashed.

\paragraph{Long and nested strings.} A string may be as long as memory
allows: a sum of 100000 terms is read in a time proportional to its length,
and takes no more of the stack of the program than a sum of two
terms.\newinvfive{} The reader of GAOL~4 built a tree with a node for each
operation, which was evaluated and deleted with one call within another per
node: a sum of 100000 terms overflowed the stack and crashed the program.
The operations are now computed as they are read, and the expressions that
a program builds through \file{gaol/parser/gaol\_expression.h} are
evaluated and deleted without recursion too. A token (a number, a name or a
run of blanks) holds $2^{29}$ characters at most, the buffer of the lexer: a
longer one throws \code{std::bad\_alloc}. The nesting of a string is limited
by the stack of the parser, which holds 10000 entries. An open parenthesis
or bracket, or a sign, takes one entry until the expression it opens is
complete; the name and the parenthesis of a call take two; an operand and
the operator or the comma after it take two, until the operand that follows
them is read. About 10000 nested parentheses, brackets or signs, 5000 nested
calls of one argument and 3300 levels of \code{1+(1+(\ldots))} are read; a
string nested deeper is refused as any string that is not an expression,
and throws \code{std::invalid\_argument}.

A \code{std::string}, given to \code{textToInterval()} or read by
\code{operator>{}>}, is read whole: a NUL character is no character of any
token, and a string holding one is not an expression.\newinvfive{} The
reader of GAOL~4 read it up to the NUL, so that \code{"[1,2]"} followed by
a NUL and any text was \itv{1}{2}, and a string of $2^{32}+3$ characters
was read as its first 3, as \code{"1+1"} for \code{"1+1+1+\ldots"}: the
lexer of flex held the length of a string in an \code{int}.

\paragraph{Errors.} A string that is not an expression throws
\code{std::invalid\_argument}. An error found anywhere in the string stops
the reading, and no interval is given: \code{"[nth\_root(8,1.5)]+[1,2]"}
throws, where the two grammars of GAOL~4 gave
\itv{-\infty}{+\infty}.\newinvfive{} Where memory runs out, the reading
throws \code{std::bad\_alloc}, and frees all it allocated.\newinvfive{} In
GAOL~4, flex wrote ``out of dynamic memory'' and ended the program, and the
parser, which read a stack it could not grow as a string nested too deep,
threw no such exception. Without exceptions, an error that the reader finds
as it computes (a wrong call, an uncertain number as a bound, two different
bounds in \code{<a, b>}, or a bound of it that is no double) makes the
reading fail, as a syntax error does: \code{gaol\_core::parse\_interval()}, of
\file{gaol/parser/gaol\_parser.h}, returns
\code{false}, and \code{gaol::textToInterval()} and \code{operator>{}>}
report the error and abort, as for any string that is not an
expression.\newinvfive{} GAOL~4 reported the error, and aborted, inside the
reader.

\begin{Ex}
The program below reads intervals from strings of the language: a literal
with a computed bound, the constants, the form of \code{textToInterval()}
with two strings, a literal divided by a literal, the bounds left out, a
hexadecimal number, the uncertain form, and a literal in a sum and in a
bound. What the printing statements write follows them.
\begin{example}
interval x = textToInterval("[4, 6*7]");
interval y = textToInterval("[-inf, dmax]");
interval z = textToInterval("[3.14,3.15]/8", "[3.14,3.15]/7");
interval t = textToInterval("[3.14,3.15]/[7,8]");
cout << textToInterval("[1,]") << ' ' << textToInterval("[ ]") << ' '
     << textToInterval("[Entire]") << endl;
@outputs^[1, inf] [empty] [-inf, inf]~
cout << textToInterval("[0x1.3p-1, 2/3]") << ' '
     << textToInterval("3.56?1") << endl;
@outputs^[0.59375, 0.6666666666666668] [3.549999999999999, 3.570000000000001]~
cout << textToInterval("1+[1,2]") << ' '
     << textToInterval("[cos([0,1]), 2]") << endl;
@outputs^[2, 3] [0.5403023058681396, 2]~
\end{example}
\end{Ex}

\section{Writing intervals}\label{sec:writing}

An interval is written to a stream with the operator \code{<{}<}, as any
primitive type of C++ is. The text depends on the active format
(Section~\ref{sec:output-format}); the empty set is written the same way in
every format.
\cindex{interval!writing}%

\begin{deffun}{ostream\&}{operator<{}<}{(ostream\& \var{out}, const interval\& \var{I})}
  Prints the interval \var{I} to the output stream \var{out}. The way the
  intervals are displayed depends on the active format
  (Section~\ref{sec:output-format}). Whatever the format, an empty interval
  is displayed as \code{[empty]}, as IEEE 1788-2015 writes the empty set
  (12.11.3).\newinvfive{} The formats \element{width} and \element{center}
  of GAOL~4 wrote \code{empty}, which is no interval literal.
  \cindex{empty!interval: printing}%
  \index[idx]{empty@\code{[empty]}}%
\end{deffun}

\noindent
In the decimal formats, each bound is rounded outward, whatever the C
library: its magnitude is written by the C library, rounded to nearest, then
compared exactly with the bound, and its last digit is moved by one when it
is on the wrong side. The text read back therefore encloses the interval
written.\newinvfive{} GAOL~4 let the C library round each bound in the
direction it set, which the C runtime of Windows and musl do not honour.

\begin{deffun}{std::string}{exact\_string}{(const interval\& \var{I})}
  Returns what \code{operator<{}<} writes in the format \element{hexa}
  (Section~\ref{sec:output-format}), without reading nor changing the global
  output format: \code{"[empty]"} for the empty set, \code{"[a]"} for a
  degenerate interval (\code{"[0x0p+0]"} for a zero of either sign), and
  otherwise the two bounds in the hexadecimal-significand form of IEEE
  1788-2015 (13.4.1), which \code{textToInterval()} reads back as the same
  doubles, a zero degenerate interval as \itv{+0}{+0}, the same set. This is
  the exact text representation of the standard, its
  \ieee{intervalToExact}.\newinvfive{} The text is the same whatever the
  locale of the program, a locale writing a decimal comma included, and
  whatever the C library.
\begin{example}
cout << exact_string(textToInterval("0.1"));
@outputs^[0x1.9999999999999p-4, 0x1.999999999999ap-4]~
\end{example}
\end{deffun}

\subsection{Converting intervals to strings}

The class \code{interval} provides a conversion operator to the standard
type \code{std::string}, which writes the interval in the active format.
\cindex{interval!conversion to \code{string}}%
\index[idx]{string@\code{string}!converting interval to|see{interval}}

\begin{example}
interval I(3,4);
string s = "test line embedding " + string(I) + " as a string";
@rem^// Now, s is "test line embedding [3, 4] as a string"~
\end{example}

\subsection{The output format}\label{sec:output-format}

Intervals may be displayed in five formats, chosen by the static method
\code{interval::format()} described after them.
\cindex{output format}%
\cindex{format!of the output}%

\begin{itemize}
\item[$\bullet$] \element{agreeing}. All the digits that are the same in
the left and in the right bound are printed once, followed by an interval
holding the remaining digits:
\begin{center}
``\element{3.141\textasciitilde[5926, 6001]}'' stands for
``\element{\itv{3.1415926}{3.1416001}}''
\end{center}
Each bound is written with the digits of the precision and the flag
\code{showpoint}, in the global locale of the program. The digits are
shared only by two bounds that are finite, not zero and of the same sign,
written with as many digits before the point and the same exponent, and
that share their first digit that is not 0: the other intervals are printed
as in the format \element{bounds}. The digits of each bound are
kept, but for the zeros that end them, \itv{1.25}{1.2567} being
\code{1.25\textasciitilde[0, 67]}.\newinvfive{} GAOL~4 dropped from both bounds what followed the last
character of the left one that is not a zero, and wrote \itv{1.25}{1.2567}
\code{1.25\textasciitilde[, ]}. It shared the digits of every interval
whose right bound is at most ten times the left one, \itv{1}{10} being
\code{1\textasciitilde[., 0.]}, \itv{4}{6} \code{\textasciitilde[4., 6.]},
without a shared digit, and a zero \code{\textasciitilde[-0., 0.]}, and
printed every interval with a negative bound with its bounds. A point
interval that the digits write exactly is printed as a single number,
without the brackets: \code{interval(4.0)} is \code{4.000000000000000},
where the format \element{bounds} prints \code{[4]}; any other point
interval is printed as any interval is, \code{interval(0.1)} being
\code{0.100000000000000\textasciitilde[0, 1]}. The point zero is printed
\code{[0]}, as in the format \element{bounds} (the formats \element{width}
and \element{center} print it \code{0}, and \element{hexa}
\code{[0x0p+0]}).
\begin{example}
interval::format(interval_format::agreeing);
std::cout << interval(4,6) << ' ' << interval(1.25,1.2567) << ' '
          << interval(-1.2567,-1.25) << std::endl;
@outputs^[4, 6] 1.25@textasciitilde[0, 67] -1.25@textasciitilde[67, 0]~
\end{example}

\item[$\bullet$] \element{bounds}. The bounds are printed between square
brackets. A degenerate interval whose value is written exactly by the
digits displayed is printed as that value alone between square brackets,
the literal of IEEE 1788-2015 for a point (12.11), which
\code{textToInterval()} reads back as the same point:
\begin{example}
interval::format(interval_format::bounds);
std::cout << interval(3,5) << ' ' << interval(4) << std::endl;
@outputs^[3, 5] [4]~
\end{example}
A zero is printed \code{[0]} whatever the signs of its bounds:
\itv{-0}{+0}, which \code{interval::zero()} is with the SSE2 and the NEON
intervals, \itv{+0}{+0} and \itv{-0}{-0} are the same set, which \code{[0]}
and \code{[-0]} are both read back as, and the text is then the same on
every build (the format \element{hexa} prints it \code{[0x0p+0]}). A
degenerate interval whose value has more digits than those displayed is
printed like any other interval, with its two bounds rounded outward:
\code{interval(0.1)} is printed \code{[0.1, 0.1000000000000001]}, which is
read back, as any other text, as an interval containing it. The bounds are
printed with a decimal point, whatever the locale of the stream, so that
the reader, whose decimal point is a point whatever the locale, reads the
text back: under a locale writing a decimal comma, it would take
\code{[-2,5]} for the interval \itv{-2}{5}, and refuses
\code{[0,5, 2,5]}.\newinvfive{} GAOL~4 printed every degenerate interval
between angles, the text of its value rounded downward and upward:
\code{interval(4)} was \code{<4, 4>}, a form of GAOL that is no literal of
the standard, and \code{interval(0.1)} \code{<0.1, 0.1000000000000001>},
which its reader refused (\code{std::invalid\_argument}), the angles
standing for one double only. The reader still takes the angles
(\pxref{sec:input-format}): a text such as \code{<4, 4>}, written by
GAOL~4, is read back.

\item[$\bullet$] \element{width}. The midpoint and the radius are printed,
\code{midpoint()} and \code{rad()}, the \ieee{mid} and \ieee{rad} of IEEE
1788-2015 (12.12.8): the
interval is contained in the midpoint plus or minus the radius, which is
the smallest number that makes it so:
\begin{example}
interval::format(interval_format::width);
std::cout << interval(4,5) << std::endl;
@outputs^4.5 (+/- 0.5)~
\end{example}
The midpoint is written rounded to nearest and the radius rounded upward,
so that a radius that is not zero is never written 0. Rounded upward, a
radius may look noisy with the default 16 digits: \itv{0.1}{0.3} is written
\code{0.2 (+/- 0.1000000000000001)}, its radius being the double nearest to
$0.1$, which is a little above $0.1$; this is intended. A degenerate
interval has no radius: it is written as its midpoint alone. An unbounded
interval has an infinite radius, its midpoint being $0$ or plus or minus
the largest double (\code{[3, inf]} is
\code{1.797693134862316e+308 (+/- inf)}). The midpoint and the radius are
written with the decimal point of the locale of the stream and its grouping
of the digits, the radius without a sign, even with the flag
\code{showpos}, and a midpoint $0$ as \code{0}.\newinvfive{} GAOL~4 wrote
$(a+b)/2$ and $(b-a)/2$ rounded to nearest, which do not contain the
interval (\itv{1}{1+2^{-52}} was \code{1 (+/- 1.11e-16)}), a radius 0 for
\itv{0}{5\cdot10^{-324}}, and \code{inf} for the midpoint of
\itv{10^{308}}{1.7\cdot10^{308}}; its manual and its header described this
number as the width.

\item[$\bullet$] \element{center}. The midpoint alone is printed, as the
format \element{width} prints it.\newinvfive

\item[$\bullet$] \element{hexa}. The left and the right bounds are printed
in the hexadecimal-significand form of IEEE 1788-2015 (13.4.1), which is
exact and which \code{textToInterval()} reads back as the same doubles,
without the round-off error of a binary-to-decimal conversion: this is the
recovery requirement of the standard (13.4), met whatever the precision,
the flags and the locale of the stream. As in the format \element{bounds},
a degenerate interval is printed as its value alone between square
brackets, \code{[0x1.8p+0]} for \code{interval(1.5)}, and a zero
\code{[0x0p+0]} whatever the signs of its bounds, which is read back as
\itv{+0}{+0}, the same set. GAOL~4 printed the sixteen hexadecimal digits
of each double, which its parser did not read.\newinvfive
\cindex{hexadecimal!form}%
\begin{example}
interval::format(interval_format::hexa);
std::cout << textToInterval("0.1") << std::endl;
@outputs^[0x1.9999999999999p-4, 0x1.999999999999ap-4]~
\end{example}
\end{itemize}

\begin{Rem}
The formats \element{bounds} and \element{hexa} are the ones read back as
an input (\pxref{sec:input-format}), whatever the locale of the stream. The
formats \element{width} and \element{center} are for the eye: the midpoint
is written with the digits of the precision, so that the text
\code{c (+/- w)} does not necessarily contain the interval, where the text
of the formats \element{bounds}, \element{agreeing} and \element{hexa} does.
A radius that took the rounding of the midpoint into account would show the
interval \code{0.1}, which is one number, as \code{0.1 (+/- 5.6e-18)}, and
would not show that an interval is narrower than the digits of its
midpoint, which is what the format is for.
\end{Rem}

\medskip
The format is chosen through the following static methods.

\defstaticmethodx{interval}{void}{format}{(interval\_format::format\_t \var{f})}
\begin{defstaticmethod}{interval}{interval\_format::format\_t}{format}{(void)}
  The first form sets the format used by the subsequent printing of
  intervals. The second form returns the format in use, a value of type
  \code{interval\_format::format\_t} (see below, and
  Section~\ref{sec:output-example} for an example of use). The format is
  global to the program, not a property of a stream: it is shared by all the
  threads, and setting it while another thread writes an interval is a data
  race, one of the common errors listed in Chapter~\ref{chap:overview}
  (\pxref{sec:common-errors}).
  \code{exact\_string()} gives the hexadecimal text without changing it.
\end{defstaticmethod}

\begin{deftype}{struct}{interval\_format}
  Structure type used to choose the output format of the intervals. It has
  five possible values of type \element{interval\_format::format\_t}:
\begin{itemize}
\item \element{interval\_format::agreeing}.
\item \element{interval\_format::bounds}.
\item \element{interval\_format::width}.
\item \element{interval\_format::center}.
\item \element{interval\_format::hexa}.
\end{itemize}
\end{deftype}

\begin{Ex}
The program below prints the same interval, \code{interval::pi()}, in the
five formats. What each statement prints is given in the comment under it.
\begin{example}
interval I(interval::pi());

interval::format(interval_format::agreeing);
cout << I << "\n";
@rem^// Prints 3.14159265358979@textasciitilde[3, 4]~
@rem^// The @textasciitilde[] part is dropped if the bounds agree on all digits~
@smallskip
interval::format(interval_format::bounds);
cout << I << "\n";
@rem^// Prints [3.141592653589793, 3.141592653589794]~
@smallskip
interval::format(interval_format::width);
cout << I << "\n";
@rem^// Prints 3.141592653589793 (+/- 4.440892098500627e-16)~
@smallskip
interval::format(interval_format::center);
cout << I << "\n";
@rem^// Prints 3.141592653589793~
@smallskip
interval::format(interval_format::hexa);
cout << I << "\n";
@rem^// Prints [0x1.921fb54442d18p+1, 0x1.921fb54442d19p+1]~
\end{example}
\end{Ex}

\subsection{The number of digits}

The number of digits written for each bound is set by the static method
\code{precision()} of the class \code{interval}. The precision of the
stream, which \code{std::setprecision} sets, is not used for intervals, and
is left as it was: the numbers written after an interval keep it. The width
of the stream (\code{std::setw}) and its adjustment apply to the whole
interval. With \code{std::internal}, the fill goes between the sign of the
midpoint and its digits in the formats \element{width} and \element{center},
and between the sign and the shared digits in the format \element{agreeing},
as it does for a number, and before the \code{[} in the other formats, as
\code{std::right} has it.\newinvfive{} GAOL~4 set the precision of the
stream to that of the intervals and left it so, and \code{std::setw} padded
the \code{[} of an interval alone.
\cindex{precision!of the output}%

\begin{defstaticmethod}{interval}{std::streamsize}{precision}{(void)}
  Returns the number of digits used for printing the bounds of intervals,
  16 by default. Section~\ref{sec:output-example} gives an example.
\end{defstaticmethod}

\begin{defstaticmethod}{interval}{std::streamsize}{precision}{(std::streamsize \var{n})}
  Sets the number of digits used for printing the bounds to \var{n}, and
  returns the number of digits previously used.
  Section~\ref{sec:output-example} gives an example.
\end{defstaticmethod}

\subsection{Example}\label{sec:output-example}

\begin{Ex}
The program below sets the precision and the format, prints $\pi$, reads
the format back, and prints $\pi$ again with sixteen digits in the format
\element{width}. Its output is shown under it.
\begin{example}
#include <iostream>
#include <gaol/gaol.h>
using namespace gaol;

using std::cout;
using std::endl;

int main(void)
{
  interval::precision(4);
  interval::format(interval_format::bounds);
  cout << interval::pi() << endl;

  if (interval::format() != interval_format::bounds) {
    cout << interval::pi() << endl;
  } else {
    interval::precision(16);
    interval::format(interval_format::width);
    cout << interval::pi() << endl;
  }
  gaol::cleanup();
  return 0;
}
@outputs^[3.141, 3.142]~
@outputs^3.141592653589793 (+/- 4.440892098500627e-16)~
\end{example}
\end{Ex}

\noindent
The first call to \code{interval::format()} is unnecessary, since the
default format is \element{interval\_format::bounds}.

The constant \code{interval::pi()} is a predefined canonical interval
containing $\pi$ (\pxref{sec:interval-constants}). Here, its bounds are two
consecutive doubles, $2^{-51}$ apart: the midpoint written is one of them,
and the radius is their distance, the smallest that has the interval within
the midpoint plus or minus it.\newinvfive{} GAOL~4 wrote half of it.

\begin{Rem}
An interval written in decimal and read back gives an interval that
encloses it, whatever the precision chosen, the bounds being rounded
outward; the fewer the digits, the wider. The format \element{hexa} gives
the same interval back.
\end{Rem}

The format \element{width} is useful whenever the number of digits
displayed is insufficient to tell whether the result is a single
floating-point number or an interval too narrow for these digits to show:
      if an interval is not a single
number, the radius displayed is different from zero, whatever its size and
the number of digits, since the radius is rounded upward. A degenerate
interval is displayed without a radius.

\section{Discussion}

This chapter described the two ways an interval crosses the border of a
program. On the way in, GAOL v5 reads every string with one grammar, in
which the literals, the numbers, the constants and the functions combine
with the operators. Its numbers are read exactly, whatever their base, their
number of digits, the locale and the C library, and a number that is not a
double becomes the tightest interval of doubles containing it. Its reader
is reentrant and computes the operations as they are read, so that several
threads may read at once and the length of a string is limited by the
memory alone. On the way out, the decimal bounds are rounded outward
whatever the C library, so that a text read back encloses the interval
written, and the format \element{hexa} gives the interval back exactly.

\medskip
The formats are trade-offs. The decimal formats \element{bounds} and
\element{agreeing} are readable and enclose the interval, at the price of a
bound moved by one unit of its last digit when the text nearest to it lies
on the wrong side; the format \element{hexa} is exact and meets the recovery
requirement of the standard; the formats \element{width} and
\element{center} are for the eye, and their text does not necessarily
contain the interval. The format and the precision are global to the
program, where the width and the adjustment belong to the stream. Reading a
long number takes a time quadratic in its number of digits.

\medskip
The next chapter turns to the namespace \code{gaol\_ieee1788}, which offers
the operations under the names of IEEE 1788-2015, the reading and the
writing of intervals included, and lists where the standard and GAOL
differ.
